\section{Limitations}\label{sec:limits}

\subsection{Threats to validity}

\emph{Internal validity.} The transient ratio is computed on a blend model that is our own definition and is not calibrated against a measured display; the absolute value of a ratio is a property of the model, and only the comparison of the methods under the same model is meaningful. The blend does not represent the temporal response of a real liquid crystal, the refresh rate, or the fluctuations of the phase modulation of the device~\cite{lizana}. The methods follow different time profiles, minimal-jerk for the proposed sequence and constant velocity for the baselines, and the transient ratio depends on the peak shift between frames; the ablation with a linear profile shows that the profile moves the result by more than the gap to the induced sequence, so the comparison of the dips mixes two effects. The pairing of the seeds is by index and the random starts are not shared, so the paired intervals are not matched-start comparisons. The seeds are \nSeedsTwo{} on the acceptance task and the sort and \longSeeds{} on the long moves, and the intervals are wide; the reading rule of the benchmark is conservative and reads some real differences of the mean as ties. The seconds were measured while other jobs loaded the machine and are not a benchmark. The implementation of the per-frame method is ours, from the description of the paper, with the weights restarting each frame; a variant that carries the weights could behave differently. The phase-change reduction of Kim et al.\ is not implemented, since the authors did not use it either.

\emph{External validity.} All results are simulated, on a scalar model, on an ideal device, with Gaussian input, point traps on bins and a lattice or the sort of a $6\times6$ array. Three tasks and two long moves do not sample the geometries that experiments use. A trap array with different weights, a three-dimensional array, a device with crosstalk, or a different SLM size could move the numbers, in either direction. The step bound is calibrated on one case and reused on the others, where it was found to hold on most seeds and not on all.

\emph{Construct validity.} The transient ratio is a proxy for atom loss: the loss itself depends on trap depth, temperature and the time spent below a level, and none of those is modelled. The phase step is the quantity the literature relates to loss, and the experiments here are far from the slips at which it was observed.

\subsection{Scope and what is not claimed}

The paper does not claim that the sequence improves on per-frame Gerchberg--Saxton with phase induction on the dip, and it measures that it does not. It does not claim a guarantee that the transient ratio stays above $0.9$ for every seed: the bound of the frame-count equation is a calibration and fails on some seeds of the long moves. It makes no claim about atom survival, about hardware, or about the computing time on a graphics processor, where the published sequences report much faster cycles~\cite{lpi}. It does not claim that the minimal-jerk profile is the right profile for the dip, and its own ablation says otherwise. It compares with two published sequences that are implemented in the engine from their descriptions, not with the authors' code, and it does not claim to reproduce their experiments, only their frame computations.

The time model is a blend of two fields and the sequences are compared at the same number of frames. A comparison at the same wall-clock time, or at the same number of iterations in total, would weigh the cold start and the iteration counts differently and could change the ranking at low frame counts. The ablation of the iteration count (Section~\ref{sec:ablation}) points in that direction: with 3 iterations per frame the sequence needs less work than the per-frame method at 5 iterations, and reaches a similar mean on the acceptance task. That comparison is not made on the long moves and the sort.
